\documentclass[11pt]{article}

\usepackage[margin=25mm]{geometry}

\usepackage{amsmath,amssymb,fontspec,unicode-math}

\IfFontExistsTF{DejaVu Serif}{\setmainfont{DejaVu Serif}}{\setmainfont{Latin Modern Roman}}

\IfFontExistsTF{Latin Modern Math}{\setmathfont{Latin Modern Math}}{\setmathfont{DejaVu Math TeX Gyre}}

\usepackage{xurl,hyperref}

\setlength{\parindent}{0pt}\setlength{\parskip}{6pt}

\begin{document}

\title{Parameter-Free Hypergraph Neural Network for Few-Shot Node Classification}

\author{ELVIA preliminary AI analysis}\date{21 October 2026}\maketitle

Preliminary assessment: Major.\par

Independent scientific review remains pending. Author adoption has not been requested.

The central least-squares approximation loses its dominant term after multiplication by labels, and overlapping hyperedges defeat exact two-hop self-information removal. Residual linearization also needs qualifications; exact symmetric prototype decisions survive.

\section*{Review scope}

Full supplied technical report read alongside the corresponding original-paper claim passages and applicable author-release source. This pass independently checks the stated reasoning and bounds; historical executions were not rerun.

\section*{Problem 1: The inverse-Gram dominant subspace is annihilated by labels}

Source: Section 3.2, Theorem 1; Appendix A.2

Assessment: supported; impact: major

\subsection*{Source claim (paraphrase)}

For balanced classes with within-class similarity \(1-\epsilon\) and between-class similarity \(\epsilon\), the paper approximates least-squares weights by \(\epsilon^{-1}Z^TY_t\).

\subsection*{Analysis}

The exact block Gram is \(M=\epsilon I+(1-2\epsilon)(I_c\otimes J_k)+\epsilon J_{kc}\). Its within-class contrast projector \(Q_1=I-(I_c\otimes J_k)/k\) has eigenvalue \(\epsilon\), but \(Q_1Y_t=0\). Therefore the supposed dominant inverse term vanishes. The surviving eigenvalues are \(\lambda_2=k(1-2\epsilon)+\epsilon\), \(\lambda_3=\lambda_2+kc\epsilon\), giving \(W_*=Z^T(\lambda_2^{-1}Q_2+\lambda_3^{-1}Q_3)Y_t\), of leading scale \(1/k\), not \(1/\epsilon\). Unit-norm features \(x_{a,j}=\sqrt\epsilon e_{a,j}+\sqrt{1-2\epsilon}f_a+\sqrt\epsilon g\) realize this exact positive-definite model, including disjoint held-out rows. Under this exact symmetry, query scores differ from prototype scores only by positive scale and a class-independent shift, so their argmax still agrees.

\subsection*{Finding and scope}

The weight-matrix approximation and its derivation fail; the practical classifier can nevertheless retain exact decisions under class symmetry. Approximate-data accuracy requires separate conditioning and margin assumptions.

\subsection*{Proposed correction}

Use the exact block-projector formula and prove decision equivalence directly. Do not extend that equivalence to unspecified approximate similarity without quantitative perturbation and margin control.

The described correction remains a proposal; author adoption or a corresponding erratum is not established.

\section*{Problem 2: Overlapping hyperedges leave unremoved second-order self-information}

Source: Section 3.3, Lemma 3, Eq. (11); Appendix Lemma 9; zen.py, RSIs

Assessment: supported; impact: major

\subsection*{Source claim (paraphrase)}

The claimed two-hop correction squares individual edge contributions, omitting cross terms when a node pair shares multiple hyperedges.

\subsection*{Analysis}

Define \(C_{ij}=\sum_{e\ni i,j}(d_e-1)^{-1}\) for distinct nodes. Matrix multiplication gives true diagonal \(r_i=d_i^{-1}\sum_{j\ne i}C_{ij}^2/(d_j-1)\), which squares the entire common-edge sum. For incidence \(H=J_4-I_4\), every node/edge degree is \(3\), every pair shares two edges, \(A_1^*=(J_4-I_4)/3\), and the true second-hop diagonal is \(1/2\). The printed and released correction subtracts \(1/4\), leaving diagonal \(1/4\), so propagation self-weight is \(\alpha_0+\alpha_2/4\). All inverses exist and no duplicate edges occur. Linear hypergraphs, where distinct edges share at most one node, avoid the omitted terms.

\subsection*{Finding and scope}

Both formula and release violate the central claim of exact independent self-information control for allowed overlapping hypergraphs. The measured classifier may still be useful.

\subsection*{Proposed correction}

Subtract the actual squared-sum diagonal, or state a linear-hypergraph restriction. The dense correction proves algebraic validity but does not preserve the claimed sparse runtime without further analysis.

The described correction remains a proposal; author adoption or a corresponding erratum is not established.

\section*{Problem 3: Initial residual layers do not collapse into a fixed propagation and one weight}

Source: Section 3.1, Table 1; Appendix B, Table 7

Assessment: supported; impact: minor

\subsection*{Source claim (paraphrase)}

The unified linearization absorbs all learned matrices into one right factor despite repeated initial-feature injections.

\subsection*{Analysis}

Take \(A=J_2/2\), \(X=(1,0)^T\), injection \(a=1/2\), scalar weights \(W_1=2,W_2=1\). The linear residual recurrence \(H_\ell=((1-a)AH_{\ell-1}+aX)W_\ell\) gives \(H_2=(1,1/2)^T\). The stated fixed propagation yields \(PX=(3/4,1/4)^T\). A common scalar would need to be both \(4/3\) and \(2\), impossible. Different injection depths carry different suffix products of learned matrices; non-residual chains or identity pre-final weights are valid special cases.

\subsection*{Finding and scope}

The unrestricted residual equivalence fails, but ZEN remains independently definable and baseline ranking changes are not inferred.

\subsection*{Proposed correction}

Retain the path-dependent matrix products, restrict the equivalence to architectures with a common factor, or label the selected propagation as a new surrogate.

The described correction remains a proposal; author adoption or a corresponding erratum is not established.

\section*{Limits}

Preliminary AI source/argument assessment; independent human scientific adjudication remains pending.\par

No author-provided or retained verification program was executed. Published empirical results were not independently reproduced.\par

\section*{Source identity}

NeurIPS 2025 official proceedings PDF (exact retained local copy inspected)

\url{https://papers.nips.cc/paper_files/paper/2025/file/65f46f5c816b2fd5a2f08ab04a8bc4a3-Paper-Conference.pdf}

Reviewed PDF SHA-256: \nolinkurl{0fc88b193f117c74ec01c017e56086f4e0c1e10da8cb468769f24742ba4152d3}

The SHA-256 identifies the exact retained PDF inspected in this rapid source/argument check. Original supplied report bytes are retained locally. No unambiguous original PDF retrieval timestamp was recovered; file modification time was not substituted.

\end{document}
